\section{Постановка задачи}
\label{sec:task}

Целью лабораторной работы является ознакомление, анализ и приобретение навыков реализации рекуррентных нейронных сетей для задачи прогнозирования (предсказания) значений числовых последовательностей.

Реализуется универсальная модель рекуррентной сети, покрывающая все шесть архитектур из вариантов задания: сеть Джордана, сеть Элмана, сеть Джордана~--~Элмана, простая рекуррентная сеть (КП), сеть с долговременной кратковременной памятью (LSTM, ДКП) и сеть с управляемыми рекуррентными блоками (GRU). Для каждого варианта скрытый слой может использовать одну из четырёх функций активации (логарифмическую, гиперболический тангенс, экспоненциальную линейную, Leaky ReLU), а также опционально линейную функцию активации эффекторов.

Требуется:
\begin{enumerate}
	\item построить модель рекуррентной сети, решающую задачу прогнозирования последовательности;
	\item построить семейства графиков зависимостей параметров модели и качества прогноза;
	\item реализовать многомерный (многошаговый) прогноз скользящим окном;
	\item дать ответы на вопросы из списка;
	\item предоставить письменный отчёт.
\end{enumerate}

Согласно требованиям методических указаний:
\begin{itemize}
	\item предусмотреть опциональное зануление значений контекстных нейронов на каждой эпохе обучения;
	\item предусмотреть опциональное указание линейной функции активации эффекторов;
	\item для сети с долговременной кратковременной памятью и управляемыми рекуррентными блоками использовать экспоненциальную линейную функцию активации эффекторных нейронов;
	\item обеспечить настраиваемые размер скользящего окна, число контекстных нейронов и нейронов на скрытом слое, а также настраиваемые начальные веса эффекторных нейронов для сетей Джордана и Джордана~--~Элмана;
	\item использовать числовые последовательности длиной не менее 10 значений;
	\item строить многомерный прогноз (не менее трёх шагов вперёд) методом скользящего окна с сохранением значений контекстных нейронов;
	\item отображать прогнозируемые значения на область выходов эффекторных нейронов (нормализация ряда).
\end{itemize}

Обучающая выборка формируется методом скользящего окна: вход~\mdash  подпоследовательность $[x_{t-w+1}, \dots, x_t]$ длины $w$, целевое значение~\mdash  следующий отсчёт $x_{t+1}$. Ряд предварительно нормализуется в диапазон $[-1, 1]$ (min--max), чтобы значения лежали в рабочей области выходов эффекторов; после прогнозирования проводится обратное преобразование.
